\FloatBarrier
\section{Cement hydration}

% cement hydration through time \cite{john.2023} -> das Paper sollte unbedingt noch eingearbeitet werden.
% TEM images of early C3S grain hydration from \cite{yan.2023}
% TEM images of CSH \cite{richardson.2004}

The main phases of \ac{OPC} are \ac{C3S}, \ac{C2S}, \ac{C3A} and \ac{C4AF}.
On contact with water these phases dissolve and new phases like \ac{CSH}, \ac{CH}, ettringite, syngenite, gypsum and \ac{AFm} form (see Figure~\ref{fig:moeser_phase_dev}) \cite{moser.2001, stark.2003}.

\begin{figure}[!htb]
    \centering
    \begin{tikzpicture}
        \begin{axis}[
				legend style={at={(1.05,.5)},anchor=west},
                legend cell align={left},
				xmin=0, xmax=28,
				ymin=0, ymax=100,
				width=0.65\linewidth,
				height=0.4\linewidth,
				grid=major,
				grid style={line width=.1pt, draw=gray!10},
				xlabel={hydration time},
                xtick={0, 11.26, 19.35, 28},
                xticklabels={0,\SI{60}{\min},\SI{24}{\hour},\SI{28}{\day}},
				ylabel={trends in proportional\\phase development},
                ylabel style={align=center},
				ytick={}
            ]

            \addplot+ [ mark=none, red, very thick ] table [x index=4, y index=5, col sep=comma, filter discard warning=false] {figures/literature/phase_dev_moeser.csv};
            \addlegendentry{\ac{CSH}};

            \addplot+ [ mark=none, blue, very thick ] table [x index=6, y index=7, col sep=comma, filter discard warning=false] {figures/literature/phase_dev_moeser.csv};
            \addlegendentry{\ac{CH}};

            \addplot+ [ mark=none, green, very thick ] table [x index=2, y index=3, col sep=comma, filter discard warning=false] {figures/literature/phase_dev_moeser.csv};
            \addlegendentry{ettringite};

            \addplot+ [ mark=none, orange, very thick ] table [x index=0, y index=1, col sep=comma, filter discard warning=false] {figures/literature/phase_dev_moeser.csv};
            \addlegendentry{syngenit};

            \addplot+ [ mark=none, grey, very thick ] table [x index=8, y index=9, col sep=comma, filter discard warning=false] {figures/literature/phase_dev_moeser.csv};
            \addlegendentry{secondary gypsum};

            \addplot+ [ mark=none, yellow, solid, very thick ] table [x index=10, y index=11, col sep=comma, filter discard warning=false] {figures/literature/phase_dev_moeser.csv};
            \addlegendentry{\ac{AFm}};


			\node[right,green] at (axis cs: 5,23)  {\small $\cong$ \SI{0.5}{\micro\meter}};
			\node[right,green] at (axis cs:21,43)  {\small $\cong$ \SI{2.5}{\micro\meter}};

			\node[right,red] at (axis cs:12.5,80)  {\small $\cong$ \SI{0.6}{\micro\meter}};
			\node[right,red] at (axis cs:21,80)  {\small $\cong$ \SI{1.0}{\micro\meter}};

        \end{axis}
    \end{tikzpicture}
    \caption{Scheme for ordinary Portland cement hydration at room temperature. The colored dimensions indicate the length of the crystal needles of \ac{CSH} and ettringite Redrawn after \textcite{moser.2001} and \textcite{stark.2003}.}
    \label{fig:moeser_phase_dev}
\end{figure}

These reaction products are responsible for the hardening of cement based materials.
The main reaction, forming \ac{CSH} and \ac{CH} can be summarized as follows \cite{gartner.2002}:
{
    \reaction{C3S + (3 - $x$ + $n$)H -> C_{x}SH_{n} + (3-$x$)CH }
    \reaction{C2S + 4.3H -> C_{1.7}SH4 + 0.3CH }
}

The hydration process can be subdivided into five periods as indicated in Figure~\ref{fig:opc_dca}.

%CEM I 42,5 R; Lengerich/Dyck.; Z808; 24.01.2023 - E-mail von Sophie vom 24.10.2024
\begin{figure}[!htb]
    \centering
    \begin{tikzpicture}
        \begin{axis}[
				legend style={at={(1.05,.5)},anchor=west},
                legend cell align={left},
				xmin=0, xmax=48,%270,
				ymin=0, ymax=30,
				width=.95\linewidth,
				height=.38\linewidth,
				grid=major,
				xlabel={Time in \si{\hour}},
				ylabel={$dQ/dt$ \si{\joule/\gram}},
				xtick={0,6,12,24,48},
				grid style={line width=.1pt, draw=gray!10}
            ]

			\addplot[mark=none, red, dashed] coordinates {(1.0,30) (1.0,0)};
			\addplot[mark=none, red, dashed] coordinates {(2.60,30) (2.60,0)};
			\addplot[mark=none, red, dashed] coordinates {(7.77,30) (7.77,0)};
			\addplot[mark=none, red, dashed] coordinates {(24,30) (24,0)};

			\node[right,preaction={fill, white, opacity=0.5}] at (axis cs:-0.1,27)  {$\leftarrow$ 0-I};
			\node[right] at (axis cs:0.75,22)  {II};
			\node[right] at (axis cs:3.5,22)  {III};
			\node[right] at (axis cs:14,22)  {IV};
			\node[right] at (axis cs:35,22)  {V};

            \addplot[blue] table [skip first n=1, x index=0, y index=1, col sep=comma] {figures/literature/DCA_OPC.csv};
        \end{axis}
    \end{tikzpicture}
    \caption{Heat flow calorimetry result of the hydration of an \ac{OPC} and its reaction periods as in \cite{gartner.2002}.}
    \label{fig:opc_dca}
\end{figure}

As described in \cite{gartner.2002}, the periods are:
\begin{multicols}{2}
  \begin{itemize}
    \item [0] Initial fast reaction
    \item [I] First deceleration period
    \item [II] Induction phase
    \item [III] Acceleration phase
    \item [IV] Second deceleration period
    \item [V] Final slow reaction
  \end{itemize}
\end{multicols}

% noch ein paar Quellen damits schön aussieht :D
The early hydration process is important to describe the initial hardening process and the rheological behavior of the cement.
The first useful qualitative observations at room temperature can be made using the \ac{SE}-mode and fresh fractured surfaces \cite{stark.2001,voland.2000,suss.2001} after a few minutes or up to \SI{10}{\hour}.
The the \ac{CSH}-needles are up to \SI{300}{\nano\meter} long, while the hydration front is barely visible in the particle \cite{voland.2000}.
% However, for (electron) microscopical analysis it is hard to get reliable data in the first few hours.
% find publications from Christiane/Long with those early hydration products
However, reliable quantitative observations on polished surfaces using the \ac{BSE} mode can fist be made after the initial hardening after a few hours.

\section{Growth of Calcium-Silicate-Hydrate}

\ac{CSH} is the decisive phase for the strength of the concrete during the later hydration process.
The detailed knowledge of the \ac{CSH}-development in volume but also microstructural shape is especially important to simulate its growth \cite{breugel.1997,nguyen-tuan.2023,nguyen-tuan.2024} and the resulting pore space.

\subsection{CSH morphology and structure}

In the literature, \ac{CSH} is usually distinguished in two different morphological types: a high density and a low density variant.
Their definition and terms vary between publications and may sometimes not mean exactly the same phenomenon:

\begin{itemize}
    \item inner / outer product \cite{scrivener.1986, richardson.1993}
    %\item phenograins and groundmass \cite{diamond.1993} % somehow this seems to be something else - read!
    \item low density (LD) / high density (HD) \ac{CSH} \cite{jennings.2000}
    \item low stiffness / high stiffness \ac{CSH} \cite{constantinides.2004}
\end{itemize}
In this document they will be called \ac{CSHi} for the dense inner structure and \ac{CSHo} for the needle or foil-like outer structure.

\ac{TEM} experiments by \textcite{melzer.1989} identified five different variations of \ac{CSH}:
\begin{itemize}
    \item gel-like \ac{CSH}: vermicular-like;surrounding anhydrous \ac{C3S} (\ac{CSHi})
    \item $\alpha$-\ac{CSH}:, type 1 needle-like, type 2 sword-like (more crystalline); projecting radially from the gel-like \ac{CSH} ((\ac{CSHo}))
    \item $\beta$-\ac{CSH}: sheet like particles; often conglomerated; also epitaxal intergrowth with $\gamma$- and $\delta$-\ac{CSH} and \ac{CH}
    \item $\gamma$-\ac{CSH}: epitaxial growths with $\beta$-\ac{CSH} and \ac{CH}
    \item $\delta$-\ac{CSH}: epitaxial growths with $\beta$- and $\gamma$-\ac{CSH} and \ac{CH}
\end{itemize}

\subsection{\ac{CSH} - gel or crystal?}

Especially in early publications, the hydration products of cement were called cement gel \cite{lea.1937,powers.1958}.
This was mainly due to to the physical properties of fresh cement paste \cite{lea.1937}.
In recent literature, the \ac{CSH}-phase is still often referred to as `gel' \cite{masoero.2018, walkley.2019, duque-redondo.2022}. %hou.2015,
According to \textsc{IUPAC}, a gel is defined as:
\begin{iquote}
    Non-fluid colloidal network or polymer network that is expanded throughout its whole volume by a fluid. \cite{gold.2019}
\end{iquote}
In the literature, two main reasons can be found for the classification of \ac{CSH} as gel:

Often, \ac{CSH} is referred to as a `gel' due to its X-ray amorphous nature compared to crystalline \ac{CSH} or tobermorite \cite{jennings.1981, garbev.2004}.
So one may argue, that water bound within the layers of the \ac{CSH} is the fluid throughout the unordered `\ac{CSH}-network' instead of bound crystal water.
However, it was shown by electron diffraction in (Cryo-)\ac{TEM} that mature \ac{CSH} phases are (crypto-)crystalline \cite{kurczyk.1960, grudemo.1960, grudemo.1965, melzer.1989, rossler.2006} with a similar structure to the tobermorite group \cite{grudemo.1960,kurczyk.1960,biagioni.2015}.
{\color{red} However, during its very early development, \ac{CSH} forms water rich, droplets, which then aggregate on top of \ac{C3S} and \ac{C2S} surfaces in an amorphous state \cite{sowoidnich.2023}. The forming amorphous first layered could then be understood as a `gel'. }
% Mail von THomas vom 04.04.2025:
% wenn man das streng nach Definition behandelt, dann sind die droplets, die wir im Nucleation…-paper gefunden haben ein Gel. War mir so noch nicht bewusst. In der Biomineral-Community sagt man eben droplet oder liquid precursor dazu. Das sind wohl die Bausteine, aus denen CSH aufgebaut wird.

% Ich würde auf jeden Fall noch sagen, dass das frühe C-S-H (young CSH in deiner Diss genannt) einen layer bildet (passivating layer) und daraus dann das ältere CSH entsteht. Hier gibt es heiße Diskussionen und Literatur, auch von uns. Im neuen paper (was wir gerade bearbeiten) werden weitere Indizien für die Existenz eines passivating layers gezeigt…(Experimente mit Freikalk und Wasserglas)
% \cite{sowoidnich.2019}
% \cite{naber.2019}
% \cite{bellmann.2015}
% \cite{bellmann.2012}
% \cite{bellmann.2010}

% Der Satz “However the composition of the young CSH is very variable in this time.” stimmt nicht so ganz: die Stöchiometrie und die Löslichkeit dieses CSH-layers sind sehr gut definiert, siehe \cite{sowoidnich.2019}

% Auch weiß ich nicht, ob die Bezeichnungen alpha-, beta-, gamma-CSH verschiedene Bedeutungen haben. In Nonat-papers werden diese Bezeichnungen auch verwendet:
% \cite{nonat.2004}
% \cite{haas.2015}

% Hier noch ein interessanter review über CSH:
% \cite{ruiz-agudo.2024}


%\textcite{grudemo.1987} compared his proposed \ac{CSH} structure to \ac{CH}, which is crystalline and can not be classified as gel.
In 1979, \textcite{ghosh.1979} argued based on morphology, that \ac{CSH} are needle-like crystals as shown in 1960 by \textcite{kurczyk.1960} and \textcite{grudemo.1960} and therefore not are not a gel as considered till then.
Also the first published \ac{SEM} images from 1966 not only show ettringite needles, but also \ac{CSH} phases in needle and platelet form are recognizable \cite{chatterji.1966}.
However, since no \ac{EDX} system was available, they only could not identify the phases without a doubt.
\textcite{stark.2001} pointed to the usually required preparation steps drying and conductive coating of the specimen for \ac{SEM} analysis until the mid-1990ies.
This preparation often leads to a significant change in the morphology of the \ac{CSH} and looking much more gel-like (e.g.: \cite{jennings.1981}).% more literature!! \cite{bergstrom.1992} has some rather bad ESEM images, however they show very young CSH and a lot of surface scattering making it look a little gel like
In a literature review \textcite{gartner.2002} also pointed to a lot of studies of morphologically damaged \ac{CSH} phases in \ac{TEM} and \ac{SEM} images, which were caused by the preparation and too aggressive imaging parameters.
The difference in preparation methods was also demonstrated by \textcite{moser.1998} using an Environmental \ac{SEM}.
\textcite{rossler.2006} were able to show, that approx. \SI{8e2}{\elementarycharge\per\square\angstrom}  at \SI{-172}{\degreeCelsius} allow to retrieve electron diffraction patterns and artefact-free imaging \ac{CSH}, while \SI{6.2e3}{\elementarycharge\per\square\angstrom} at room temperature induce significant beam damage to the delicate crystals.
Also, at very high magnifications, the energy dose is very high even at low voltages and currents.
This can cause \ac{CSH}-structures to bent or melt and introduce pores within originally compact phases (Figure~\ref{fig:ebeam_damage}) \cite{kleiner.2021a,kleiner.2021f}.

\begin{figure}[h!]
    \centering

	\ifdefined\PRINTABLE
		\includegraphics[height=0.33\textheight]{methods/ebeam_damage.pdf}
	\else
		\animategraphics[palindrome,autoplay,height=0.33\textheight]{3}{methods/ebeam_damage_anim/Animation}{0}{4}
	\fi

    \caption{
        Electron beam damage of a \ac{CSH} needle of \SI{28}{\day} hydrated alite.
        \ac{SE} images captured at \SI{350}{\volt} accelerating voltage and \SI{200}{\volt} stage bias.
        The image set was captured within \SI{124}{\second} during continuous electron beam exposure \cite{kleiner.2021a,kleiner.2021f}.
    }
    \label{fig:ebeam_damage}
\end{figure}

Furthermore, early \ac{SEM}s required higher voltages to image the surfaces, causing subsurface scattering and reduced surface detail \cite{moser.1998}, increasing the impression of a gel like material due to reduces surface details.

However, a gel can also be defined as a substance capable of swelling in the presence of a second component (in this case water) \cite{hermans.1946}.
This seems to be the case, as observable by the creep, swelling and shrinkage behavior of concrete \cite{powers.1958,feldman.1970}, which can be attributed to \ac{CSH} \cite{zhou.2025}.

The colloidal nature of hydration products of cement was proposed more than a century ago, as a lively discussion between \textcite{michaelis.1908} and the editorial of the \textsc{Tonindustrie-Zeitung} from 1907/1908 demonstrates.
The required colloidal nature (porous structure between \SI{1}{\nano\meter} to \SI{1}{\micro\meter}) can be applied to water saturated, low density \ac{CSHo} and very young \ac{CSH} \cite{sowoidnich.2023}.
However, according to \ac{SEM}-images dense \ac{CSHi} has a very low porosity and its morphology can hardly be classified as colloidal.
Therefore, the swelling is often be attributed to the interlayer porosity of \ac{CSH} \cite{jennings.2008}, which is usually in the range of \num{0.9} to \SI{1.4}{\nano\meter} for tobermorite-like structures, while the latter seems to be more fitting for \ac{CSH} as discussed by \textcite{duque-redondo.2022}.
At this pore size, the colloidal criteria for a gel would be fulfilled.

In 1937, \textcite{lea.1937} called hardened cement a `rigid gel' and described similarities to silica gel, which is today classified as a `xerogel' (a solid dried gel).
Accordingly, \textcite{wittmann.1976,wittmann.1977} described \ac{CSH} phases as xerogel in the 1970ies.
However, some authors like \textcite{almdal.1993} argue, that `\textit{xerogel and undiluted systems should not be admitted as gels}'.
Furthermore, they discussed, that a gel should contain a `\textit{substantial quantity}' of liquid \cite{almdal.1993}, which is debatable for \ac{CSH}, since a lot of water is bound in its crystalline structure.

In 2001 \textcite{stark.2001} concluded, that while the designation \ac{CSH}-gel is correct `\textit{from a physical point of view as gel implies a particle size}', `\textit{from the point of view of cement chemistry and mineralogy}' \ac{CSH} can not be classified as gel due to its nano-crystalline nature.

Considering the literature, the term `\ac{CSH}-gel' will not be used in this work due to the nano-crystalline nature of \ac{CSH}.
While \ac{CSH} forms a colloidal network, it forms a rigid solid, not the soft, elastic material often associated with the common understanding of a gelatin like `gel'.

\subsection{\ac{CSH} layer development}

In 1929 / 1930 \textsc{Anderegg} and \textsc{Hubbell} \cite{anderegg.1929,anderegg.1930} already tried to measure the hydration layer development of \ac{OPC} at \SI{21(3)}{\celsius} using optical light microscopy.
They utilized the difference of the refraction index between \ac{CSH} and unhydrated clinker.
According to their experiments, the hydrate rim is depending on the mean particle dimension (measured from 150 particles) as shown in Figure~\ref{fig:anderegg-hubbel-distribution}.

\begin{figure}[!htb]
    \centering
    \begin{tikzpicture}
        \begin{axis}[
				legend style={at={(1.05,.5)},anchor=west},
                legend cell align={left},
				xmin=10, xmax=30,%270,
				ymin=0, ymax=7,
				width=0.7\linewidth,
				height=0.41\linewidth,
				grid=major,
				xlabel={mean particle dimension in \si{\micro\meter}},
				ylabel={$l_{\text{CSH}}$ in \si{\micro\meter}},
				ytick={0,1,...,7},
				grid style={line width=.1pt, draw=gray!10}
            ]

            \addplot+ [ only marks, mark=x, orange ] table [skip first n=1, x index=0, y index=3, col sep=comma, filter discard warning=false] {figures/hydrate-rim/data_andregg/gsd-PCA.csv};
            \addplot [ draw=none, forget plot ] table [skip first n=1, col sep=comma, x index=0, y={create col/linear regression={y=[index]3}}, filter discard warning=false ] {figures/hydrate-rim/data_andregg/gsd-PCA.csv};
            \addplot[update limits=false, no markers, orange, domain=15:26, samples=2, forget plot] {\pgfplotstableregressionb + x*\pgfplotstableregressiona};
            \addlegendentry{\SI{7}{\day}, \acs*{OPC}};

            \addplot+ [ only marks, mark=o, orange ] table [skip first n=1, x index=0, y index=3, col sep=comma, filter discard warning=false] {figures/hydrate-rim/data_andregg/gsd-PCA+CaSO4.csv};
            \addplot [ draw=none, forget plot ] table [skip first n=1, col sep=comma, x index=0, y={create col/linear regression={y=[index]3}}, filter discard warning=false ] {figures/hydrate-rim/data_andregg/gsd-PCA+CaSO4.csv};
            \addplot[update limits=false, no markers, orange, dashed, domain=16:26, samples=2, forget plot] {\pgfplotstableregressionb + x*\pgfplotstableregressiona};
            \addlegendentry{\SI{7}{\day}, \acs*{OPC}+\ce{CaSO4}};

            %\addplot[orange, domain=14:30, opacity=0.5, samples=2] {0.068516925009012*x + 1.05223737455958};
            %\addlegendentry{\SI{7}{\day}, $R^2=0.11$};

            \addplot+ [ only marks, mark=x, cyan ] table [skip first n=1, x index=0, y index=4, col sep=comma, filter discard warning=false] {figures/hydrate-rim/data_andregg/gsd-PCA.csv};
            \addplot [ draw=none, forget plot ] table [skip first n=1, col sep=comma, x index=0, y={create col/linear regression={y=[index]4}}, filter discard warning=false ] {figures/hydrate-rim/data_andregg/gsd-PCA.csv};
            \addplot[update limits=false, no markers, cyan, domain=15:28, samples=2, forget plot] {\pgfplotstableregressionb + x*\pgfplotstableregressiona};
            \addlegendentry{\SI{28}{\day}, \acs*{OPC}};

            \addplot+ [ only marks, mark=o, cyan ] table [skip first n=1, x index=0, y index=4, col sep=comma, filter discard warning=false] {figures/hydrate-rim/data_andregg/gsd-PCA+CaSO4.csv};
            \addplot [ draw=none, forget plot ] table [skip first n=1, col sep=comma, x index=0, y={create col/linear regression={y=[index]4}}, filter discard warning=false ] {figures/hydrate-rim/data_andregg/gsd-PCA+CaSO4.csv};
            \addplot[update limits=false, no markers, cyan, dashed, domain=16:26, samples=2, forget plot] {\pgfplotstableregressionb + x*\pgfplotstableregressiona};
            \addlegendentry{\SI{28}{\day}, \acs*{OPC}+\ce{CaSO4}};

            %\addplot[cyan, domain=14:30, opacity=0.5, samples=2] {0.196884840055179*x + 0.0780310618817959};
            %\addlegendentry{\SI{28}{\day}, $R^2=0.59$};

            \addplot+ [ only marks, mark=x, red ] table [skip first n=1, x index=0, y index=5, col sep=comma, filter discard warning=false] {figures/hydrate-rim/data_andregg/gsd-PCA.csv};
            \addplot [ draw=none, forget plot ] table [skip first n=1, col sep=comma, x index=0, y={create col/linear regression={y=[index]5}}, filter discard warning=false ] {figures/hydrate-rim/data_andregg/gsd-PCA.csv};
            \addplot[update limits=false, no markers, red, domain=15:28, samples=2, forget plot] {\pgfplotstableregressionb + x*\pgfplotstableregressiona};
            \addlegendentry{\SI{90}{\day}, \acs*{OPC}};

            \addplot+ [ only marks, mark=o, red ] table [skip first n=1, x index=0, y index=5, col sep=comma, filter discard warning=false] {figures/hydrate-rim/data_andregg/gsd-PCA+CaSO4.csv};
            \addplot [ draw=none, forget plot ] table [skip first n=1, col sep=comma, x index=0, y={create col/linear regression={y=[index]5}}, filter discard warning=false ] {figures/hydrate-rim/data_andregg/gsd-PCA+CaSO4.csv};
            \addplot[update limits=false, no markers, red, dashed, domain=16:26, samples=2, forget plot] {\pgfplotstableregressionb + x*\pgfplotstableregressiona};
            \addlegendentry{\SI{90}{\day}, \acs*{OPC}+\ce{CaSO4}};

            %\addplot[red, domain=14:30, opacity=0.5, samples=2] {0.224971630191457*x + 0.566670110374461};
            %\addlegendentry{\SI{90}{\day}, $R^2=0.82$};
        \end{axis}
    \end{tikzpicture}
    \caption{Mean hydration rim of \ac{OPC} from experiments of \textsc{Anderegg} and \textsc{Hubbell} \cite{anderegg.1929,anderegg.1930} in dependence of the mean particle `dimension' of 150 measured grains. The lack or addition of \SI{0.8}{\percent} \ce{CaSO4} in the experiments is marked with $\times$ and $\circ$ respectively. }
    \label{fig:anderegg-hubbel-distribution}
\end{figure}

This dependence becomes more prone for older specimens, since particles with a radius equal or smaller than the hydrate layer are consumed and cannot contribute to a larger hydrate layer.
However, they mainly observed rather coarse particles.

In 1930 they also published datasets for the hydration of \ac{C3S} and \ac{C2S} \cite{anderegg.1930}.
In contrast to the \ac{OPC} experiments, these consisted only out of very few series of experiments with a rather high variance (see Figure~\ref{fig:anderegg-hubbel}).

\textcite{garrault.2006} measured the development of the \ac{CSH} concentration of \SI{4}{\gram} \ac{C3S} grains with different median diameters (7, 10, 11, 12.5, \SI{14}{\micro\meter}) for up to one day in a concentrated lime solution.
They estimated the thickness of \ac{CSHi} (replacing \ac{C3S}) and \ac{CSHo} (on top of the grain) at the time, when the \ac{CSH} precipitation became limited by diffusion (see Figure~\ref{fig:garrault-distribution}) (beginning of the second deceleration period in Figure~\ref{fig:opc_dca}).
For \SI{7}{\micro\meter} the diffusion was limited was after $\approx$ \SI{600}{\min} (\SI{10}{\hour}), for \SI{14}{\micro\meter} after $\approx$ \SI{700}{\min} (\SI{12}{\hour}).
The estimation is based on the \ac{CSH}-concentration, the mean diameter and that the particles are spherical.
\textcite{garrault.2006} included a measurement of a single \ac{TEM}-image from  of a hydrated \ac{C3S} grain with a diameter of \SI{3}{\micro\meter}.% made by \textsc{Lachowski} and \textsc{Courault} in 1998.

\begin{figure}[!htb]
    \centering
    \begin{tikzpicture}
        \begin{axis}[
				legend style={at={(.97,.97)},anchor=north east},
                legend cell align={left},
                legend columns = -1,
				xmin=0, xmax=20,
				ymin=0, ymax=0.7,
				width=0.7\linewidth,
				height=0.35\linewidth,
				grid=major,
				xlabel={mean particle diameter in \si{\micro\meter}},
				ylabel={$l_{\text{CSH}}$ in \si{\micro\meter}},
				ytick={0,0.1,...,0.7},
				grid style={line width=.1pt, draw=gray!10},
            ]

            \addplot [ mark=+, blue, only marks ] table {figures/hydrate-rim/data_garrault/CSH.dat};
            \addplot [ draw=none, forget plot ] table [ y={create col/linear regression={y=Y}} ] {figures/hydrate-rim/data_garrault/CSH.dat};
            \addplot[update limits=false, no markers, blue, domain={2:16}, forget plot] {\pgfplotstableregressionb + x*\pgfplotstableregressiona};
            \addlegendentry{\acs*{CSH}};

            \addplot [ mark=x, red, only marks ] table {figures/hydrate-rim/data_garrault/CSHi.dat};
            \addplot [ draw=none, forget plot ] table [ y={create col/linear regression={y=Y}} ]  {figures/hydrate-rim/data_garrault/CSHi.dat};
            \addplot[update limits=false, no markers, red, domain={6:16}, forget plot] {\pgfplotstableregressionb + x*\pgfplotstableregressiona};
            \addlegendentry{\acs*{CSHi}};

            \addplot [ mark=x, orange, only marks ] table {figures/hydrate-rim/data_garrault/CSHo.dat};
            \addplot [ draw=none, forget plot ] table [ y={create col/linear regression={y=Y}} ]  {figures/hydrate-rim/data_garrault/CSHo.dat};
            \addplot[update limits=false, no markers, orange, domain={6:16}, forget plot] {\pgfplotstableregressionb + x*\pgfplotstableregressiona};
            \addlegendentry{\acs*{CSHo}};

            \addplot [ mark=diamond, black, only marks ] coordinates { (3, 0.50) };
            \addlegendentry{\acs*{TEM}};

        \end{axis}
    \end{tikzpicture}
    \caption{Mean hydration rim of \SI{4}{\gram} \ac{C3S} from an estimation based of experiments by \textcite{garrault.2006} in dependence of the mean particle diameter after roughly 10 to \SI{12}{\hour}. The \ac{CSHo} data ($l_{\text{CSHo}}=l_{\text{CSH}}-l_{\text{CSHi}}$) were calculated for better comparability.}
    \label{fig:garrault-distribution}
\end{figure}

The data of \textcite{garrault.2006} indicate a smaller hydrate seam with increasing mean particle diameters.
This is contradicting the data of \textsc{Anderegg} and \textsc{Hubbell} \cite{anderegg.1929,anderegg.1930} (see Figure~\ref{fig:anderegg-hubbel-distribution}).
However, \textsc{Anderegg} and \textsc{Hubbell} began their measurements after 7 days and therefore after the diffusion limited growth of the \ac{CSH}.
Furthermore, their mean particle diameter was significantly larger (starting from \SI{15.5}{\micro\meter})

\ifdefined\EXTERNALIZE
    \tikzexternaldisable
\fi

\begin{figure}[!htb]
    \centering
    \begin{tikzpicture}
        \begin{groupplot}[
            group style={group size=2 by 1,
                    group name=coupeplot,
                    xticklabels at=edge bottom,
                    horizontal sep=0pt},
            legend style={at={(.05,.94)},anchor=north west},
            axis x line=box,
            ymin=0, ymax=8.5,
            ytick={0,1,...,8},
            ylabel={$l_{\text{CSH}}$ in \si{\micro\meter}},
            xlabel={time $t$ in \si{\day}},
            width=0.5\linewidth,
            height=0.4\linewidth,
            legend columns = 4,
            legend style = {anchor=north, style={column sep=3pt}},
            grid=both,
            grid style={line width=.1pt, draw=gray!10},
        ]
        \nextgroupplot[
            xmin=0, xmax=29,
            restrict x to domain=0:42,
            xtick={1,9,21,28},
            axis y line*=left,
            legend to name={CommonLegend3}
        ]

            %C3S
            \addplot [ mark=x, red, opacity=0.5 ] coordinates {
                        (0,0)
                        (0.125,1.68)
                        (1,2.25)
                        %(2,3.02)
                        %(3,3.60)
                        (7,4.32)
                        (28,4.44)
                        (165,5.77)
                    };
            \addlegendentry[]{\acs*{C3S} \cite{anderegg.1930}}


            %C2S
            \addplot [ mark=x, orange, opacity=0.5 ] coordinates {
                        (0,0)
                        (1,0.28)
                        %(2,0.40)
                        %(3,0.50)
                        (7,0.62)
                        (28,0.83)
                        (165,3.54)
                    };
            \addlegendentry[]{\acs*{C2S} \cite{anderegg.1930}}

            %PC A average
            \addplot+ [draw=none, no markers, name path=A,forget plot] table [skip first n=2, x index=0, col sep=comma,y expr={\thisrowno{12}-\thisrowno{13}}, filter discard warning=false ] {figures/hydrate-rim/data_andregg/PCA.csv};
            \addplot+ [draw=none, no markers, name path=B,forget plot] table [skip first n=2, x index=0, col sep=comma,y expr={\thisrowno{12}+\thisrowno{13}}, filter discard warning=false ] {figures/hydrate-rim/data_andregg/PCA.csv};

            \addplot [fill=cyan, fill opacity=0.2,forget plot] fill between [of=A and B];

            \addplot [ mark=x, cyan, opacity=0.5 ] table [skip first n=2, x index=0, y index=12, col sep=comma, filter discard warning=false] {figures/hydrate-rim/data_andregg/PCA.csv};
            \addlegendentry[]{\acs*{OPC} \cite{anderegg.1929,anderegg.1930}}


            % \addplot+ [draw=none, no markers, name path=A,forget plot] table [skip first n=2, x index=0, col sep=comma,y expr={\thisrowno{9}-\thisrowno{10}}, ] {figures/hydrate-rim/data_andregg/PCA+CaSO4.csv};
            % \addplot+ [draw=none, no markers, name path=B,forget plot] table [skip first n=2, x index=0, col sep=comma,y expr={\thisrowno{9}+\thisrowno{10}}, ] {figures/hydrate-rim/data_andregg/PCA+CaSO4.csv};

            % \addplot [fill=green, fill opacity=0.2,forget plot] fill between [of=A and B];
            % \addplot [ mark=x, green, opacity=0.5 ] table [skip first n=2, x index=0, y index=9, col sep=comma] {figures/hydrate-rim/data_andregg/PCA+CaSO4.csv};
            % \addlegendentry[]{\acs{OPC} + CaSO4 \cite{anderegg.1929,anderegg.1930}}

            \addplot+ [
                only marks,
                mark=x, black, opacity=0.5,
                error bars/.cd,
                    y dir=both,y explicit,
            ] coordinates {
                        %(0,0) +- (0,0)
                        (1,.933) +- (0,0.083)
                        (2,1.188) +- (0,0.049)
            };
            \addlegendentry[]{\acs*{C3S} \cite{nicoleau.2011}}


        \nextgroupplot[
            xmin=0, xmax=270,
            restrict x to domain=27:1000,
            xtick={28, 90, 165, 270},
            %axis y line=right,
            axis y line*=right,
            axis x discontinuity=parallel,
        ]

            %C3S
            \addplot [ mark=x, red, opacity=0.5 ] coordinates {
                        (0,0)
                        (0.125,1.68)
                        (1,2.25)
                        %(2,3.02)
                        %(3,3.60)
                        (7,4.32)
                        (28,4.44)
                        (165,5.77)
                    };

            %C2S
            \addplot [ mark=x, orange, opacity=0.5 ] coordinates {
                        (0,0)
                        (1,0.28)
                        %(2,0.40)
                        %(3,0.50)
                        (7,0.62)
                        (28,0.83)
                        (165,3.54)
                    };

            % %C3A
            % \addplot [ mark=x, blue, opacity=0.5 ] coordinates {
            %             (0,0)
            %             (0.01042,2.14)
            %             (0.125,4.35)
            %             %(1,4.7)
            %             %(2,5.2)
            %             (3,5.68) %(7,5.70)
            %             (165,5.66)
            %         };
            % \addlegendentry[]{\acs*{C3A} \cite{anderegg.1930}}

            %PC A average
            \addplot+ [draw=none, no markers, name path=A,forget plot] table [skip first n=3, x index=0, col sep=comma,y expr={\thisrowno{12}-\thisrowno{13}}, filter discard warning=false ] {figures/hydrate-rim/data_andregg/PCA.csv};
            \addplot+ [draw=none, no markers, name path=B,forget plot] table [skip first n=3, x index=0, col sep=comma,y expr={\thisrowno{12}+\thisrowno{13}}, filter discard warning=false ] {figures/hydrate-rim/data_andregg/PCA.csv};

            \addplot [fill=cyan, fill opacity=0.2,forget plot] fill between [of=A and B];

            \addplot [ mark=x, cyan, opacity=0.5 ] table [skip first n=3, x index=0, y index=12, col sep=comma, filter discard warning=false] {figures/hydrate-rim/data_andregg/PCA.csv};


            % \addplot+ [draw=none, no markers, name path=A,forget plot] table [skip first n=2, x index=0, col sep=comma,y expr={\thisrowno{9}-\thisrowno{10}}, ] {figures/hydrate-rim/data_andregg/PCA+CaSO4.csv};
            % \addplot+ [draw=none, no markers, name path=B,forget plot] table [skip first n=2, x index=0, col sep=comma,y expr={\thisrowno{9}+\thisrowno{10}}, ] {figures/hydrate-rim/data_andregg/PCA+CaSO4.csv};

            % \addplot [fill=green, fill opacity=0.2,forget plot] fill between [of=A and B];
            % \addplot [ mark=x, green, opacity=0.5 ] table [skip first n=2, x index=0, y index=9, col sep=comma] {figures/hydrate-rim/data_andregg/PCA+CaSO4.csv};

        \end{groupplot}
        \node [above] (L) at (rel axis cs:1.0,1.04) {\ref{CommonLegend3}};

    \end{tikzpicture}
    \caption{Growth of \ac{CSH} on \ac{OPC}, \ac{C3S} and \ac{C2S} measured by \textsc{Anderegg} and \textsc{Hubbell} \cite{anderegg.1929,anderegg.1930} and \textcite{nicoleau.2011}}
    \label{fig:anderegg-hubbel}
\end{figure}

In the 1960ies \textcite{kurczyk.1960} used an electron microscope to measure the length of \ac{CSH}-needles of one month hydrated \ac{C3S}.
They measured an average length of \SI{500}{\nano\meter}.
However, the hydrated material was ground and dispersed using ultrasonic treatment, which may have damaged the hydrates.

\textcite{nicoleau.2011} measured the \ac{CSH}-thickness using Cryo-\ac{SEM} on cracked surfaces of hydrated \ac{C3S} ($w/b=0.5$, $\vartheta=$\SI{20}{\celsius}) (see Figure~\ref{fig:anderegg-hubbel}).
\textcite{bishnoi.2012} criticized this approach due to the rather uncontrolled cracking, the subjectivity of the measurement and the occurrence of the sectioning effect, which was described in detail by \textcite{scrivener.2004}.
This effect is an issue for all 2D images of polished sections (e.g. typical \ac{SEM}-\ac{BSE} images) and will be described and discussed in Subsection~\ref{subsec:sectioning-effect}.

\textcite{voland.2000} and \textcite{suss.2001} manually measured the development of the length of \ac{CSH}-needles, (roughly equivalent with \ac{CSHo}), using \ac{SEM}-\ac{SE}-images until a specimen age of \SI{28}{\day}.
\textsc{Stark} and \textsc{Möser} et al. \cite{stark.2001,moser.2001,moser.2002} continued this work close to the three years mark.
The experiments were done using \ac{C3S} and $\beta$-\ac{C2S} at different temperatures.
They measured a minimum and maximum value, of which the mean was calculated.
However, it was not made clear by the authors, how many measurements were taken for each data point. Therefore, it should be assumed that these data are based on individual measurements.
The results are plotted in Figures~\ref{fig:suess_CSH_C3S} (\ac{C3S}) and \ref{fig:suess_CSH_C2S} ($\beta$-\ac{C2S}).

\begin{figure}[!htb]
    \centering
    \begin{tikzpicture}
        \begin{groupplot}[
            group style={group size=2 by 1,
                    group name=coupeplot,
                    xticklabels at=edge bottom,
                    horizontal sep=0pt},
            legend style={at={(.05,.94)},anchor=north west},
            ylabel = {$l_{\text{CSH}_\text{o}}$ in \si{\micro\meter}},
            width=0.5\linewidth,
            height=0.4\linewidth,
            legend columns = 4,
            axis x line = box,
            legend style = {anchor=north, style={column sep=3pt}},
            grid=both,
            grid style={line width=.1pt, draw=gray!10},
            y tick label style={
                /pgf/number format/fixed zerofill,
                /pgf/number format/precision=1},,
            scaled ticks=false, tick label style={/pgf/number format/fixed}
        ]
        \nextgroupplot[
            xmin=0, xmax=635,
            ymin=0, ymax=1.250,
            xlabel={time $t$ in \si{\min}},
            xtick={0,0,130,300,600},
            ytick={0,.250,...,3.000},
            restrict x to domain=0:600,
            axis y line*=left,
            y tick label style={
                /pgf/number format/fixed zerofill,
                /pgf/number format/precision=2},
        ]
            \addplot+ [draw=none, no markers, name path=A] table [skip first n=3, x index=0, y index=2, col sep=comma, filter discard warning=false] {figures/hydrate-rim/data_t_suess/CSH-length_C3S-C2S.csv};
            \addplot+ [draw=none, no markers, name path=B] table [skip first n=3, x index=0, y index=4, col sep=comma, filter discard warning=false] {figures/hydrate-rim/data_t_suess/CSH-length_C3S-C2S.csv};
            \addplot [fill=green, fill opacity=0.2] fill between [of=A and B];

            \addplot [color=green, mark=+, dotted] table [skip first n=3, x index=0, y index=3, col sep=comma, filter discard warning=false] {figures/hydrate-rim/data_t_suess/CSH-length_C3S-C2S.csv};


            \addplot+ [draw=none, no markers, name path=A] table [skip first n=3, x index=0, y index=5, col sep=comma, filter discard warning=false] {figures/hydrate-rim/data_t_suess/CSH-length_C3S-C2S.csv};
            \addplot+ [draw=none, no markers, name path=B] table [skip first n=3, x index=0, y index=7, col sep=comma, filter discard warning=false] {figures/hydrate-rim/data_t_suess/CSH-length_C3S-C2S.csv};
            \addplot [fill=cyan, fill opacity=0.2] fill between [of=A and B];

            \addplot [color=cyan, mark=x, dotted] table [skip first n=3, x index=0, y index=6, col sep=comma, filter discard warning=false] {figures/hydrate-rim/data_t_suess/CSH-length_C3S-C2S.csv};


            \addplot+ [draw=none, no markers, name path=A] table [skip first n=3, x index=0, y index=8, col sep=comma, filter discard warning=false] {figures/hydrate-rim/data_t_suess/CSH-length_C3S-C2S.csv};
            \addplot+ [draw=none, no markers, name path=B] table [skip first n=3, x index=0, y index=10, col sep=comma, filter discard warning=false] {figures/hydrate-rim/data_t_suess/CSH-length_C3S-C2S.csv};
            \addplot [fill=red, fill opacity=0.2] fill between [of=A and B];

            \addplot [color=red, mark=star, dotted] table [skip first n=3, x index=0, y index=9, col sep=comma, filter discard warning=false] {figures/hydrate-rim/data_t_suess/CSH-length_C3S-C2S.csv};


            \addplot+ [draw=none, no markers, name path=A] table [skip first n=3, x index=0, y index=11, col sep=comma, filter discard warning=false] {figures/hydrate-rim/data_t_suess/CSH-length_C3S-C2S.csv};
            \addplot+ [draw=none, no markers, name path=B] table [skip first n=3, x index=0, y index=13, col sep=comma, filter discard warning=false] {figures/hydrate-rim/data_t_suess/CSH-length_C3S-C2S.csv};
            \addplot [fill=orange, fill opacity=0.2] fill between [of=A and B];

            \addplot [color=orange, mark=Mercedes star, dotted] table [skip first n=3, x index=0, y index=12, col sep=comma, filter discard warning=false] {figures/hydrate-rim/data_t_suess/CSH-length_C3S-C2S.csv};

        \nextgroupplot[
            xmin=-1.5, xmax=28,
            ymin=0, ymax=3.000,
            xlabel={time $t$ in \si{\day}},
            xtick={1,9,21,28},
            ytick={0.5,1,1.500,...,3.000},
            restrict x to domain=1:28,
            axis y line*=right,
            axis x discontinuity=parallel,
            legend to name={CommonLegend}
        ]
            \pgfplotsset{cycle list shift=+1}
            \addplot+ [draw=none, no markers, name path=A,forget plot] table [skip first n=3, x index=1, y index=2, col sep=comma, filter discard warning=false] {figures/hydrate-rim/data_t_suess/CSH-length_C3S-C2S.csv};
            \addplot+ [draw=none, no markers, name path=B,forget plot] table [skip first n=3, x index=1, y index=4, col sep=comma, filter discard warning=false] {figures/hydrate-rim/data_t_suess/CSH-length_C3S-C2S.csv};
            \addplot [fill=green, fill opacity=0.2,forget plot] fill between [of=A and B];

            \addplot [color=green, mark=+, dotted] table [skip first n=3, x index=1, y index=3, col sep=comma, filter discard warning=false] {figures/hydrate-rim/data_t_suess/CSH-length_C3S-C2S.csv};
            \addlegendentry[]{\ac{C3S} \SI{5}{\celsius} \cite{suss.2001}}


            \addplot+ [draw=none, no markers, name path=A,forget plot] table [skip first n=3, x index=1, y index=5, col sep=comma, filter discard warning=false] {figures/hydrate-rim/data_t_suess/CSH-length_C3S-C2S.csv};
            \addplot+ [draw=none, no markers, name path=B,forget plot] table [skip first n=3, x index=1, y index=7, col sep=comma, filter discard warning=false] {figures/hydrate-rim/data_t_suess/CSH-length_C3S-C2S.csv};
            \addplot [fill=cyan, fill opacity=0.2,forget plot] fill between [of=A and B];

            \addplot [color=cyan, mark=x, dotted] table [skip first n=3, x index=1, y index=6, col sep=comma, filter discard warning=false] {figures/hydrate-rim/data_t_suess/CSH-length_C3S-C2S.csv};
            \addlegendentry[]{\ac{C3S} \SI{20}{\celsius} \cite{voland.2000}}


            \addplot+ [draw=none, no markers, name path=A,forget plot] table [skip first n=3, x index=1, y index=8, col sep=comma, filter discard warning=false] {figures/hydrate-rim/data_t_suess/CSH-length_C3S-C2S.csv};
            \addplot+ [draw=none, no markers, name path=B,forget plot] table [skip first n=3, x index=1, y index=10, col sep=comma, filter discard warning=false] {figures/hydrate-rim/data_t_suess/CSH-length_C3S-C2S.csv};
            \addplot [fill=red, fill opacity=0.2,forget plot] fill between [of=A and B];

            \addplot [color=red, mark=star, dotted] table [skip first n=3, x index=1, y index=9, col sep=comma, filter discard warning=false] {figures/hydrate-rim/data_t_suess/CSH-length_C3S-C2S.csv};
            \addlegendentry[]{\ac{C3S} \SI{60}{\celsius} \cite{suss.2001}}


            \addplot+ [draw=none, no markers, name path=A,forget plot] table [skip first n=3, x index=1, y index=11, col sep=comma, filter discard warning=false] {figures/hydrate-rim/data_t_suess/CSH-length_C3S-C2S.csv};
            \addplot+ [draw=none, no markers, name path=B,forget plot] table [skip first n=3, x index=1, y index=13, col sep=comma, filter discard warning=false] {figures/hydrate-rim/data_t_suess/CSH-length_C3S-C2S.csv};
            \addplot [fill=orange, fill opacity=0.2,forget plot] fill between [of=A and B];

            \addplot [color=orange, mark=Mercedes star, dotted] table [skip first n=3, x index=1, y index=12, col sep=comma, filter discard warning=false] {figures/hydrate-rim/data_t_suess/CSH-length_C3S-C2S.csv};
            \addlegendentry[]{\ac{C3S} \SI{90}{\celsius} \cite{suss.2001}}

        \end{groupplot}
        \node [above] (L) at (rel axis cs:1.0,1.05) {\ref{CommonLegend}};

    \end{tikzpicture}
    \caption{Growth of \ac{CSHo} on \ac{C3S} at different temperatures. Based on data from \textcite{voland.2000} and \textcite{suss.2001}.}
    \label{fig:suess_CSH_C3S}
\end{figure}

\ifdefined\EXTERNALIZE
    \tikzexternalenable
\fi

% TODO: include data from \cite{shirani.2023}?
% - 450 - 550 nm shell thickness of csh after 19h (i do not really get why there exist these two numbers? Durchschnitt insgesamt vs hadley grain-shell?)

The formation of \ac{CSH} begins earlier at increasing temperatures.
This also causes a significant increase of the overall fibre length.
At the same temperature, \ac{C2S} demonstrates a slower \ac{CSH}-growth.
While the \ac{CSH}-growth of \ac{C3S} stagnates after 5 months, \textcite{moser.2002} indicates that the needle length of \ac{CSH} on \ac{C2S} is still increasing significantly, even until an age of 3 years.

\begin{figure}[!htb]
    \centering
    \begin{tikzpicture}
        \begin{groupplot}[
            group style={group size=2 by 1,
                    group name=coupeplot,
                    xticklabels at=edge bottom,
                    horizontal sep=0pt},
            legend style={at={(.05,.94)},anchor=north west},
            axis x line=box,
            ytick={0,1,1.5,...,3},
            ylabel={$l_{\text{CSH}_\text{o}}$ in \si{\micro\meter}},
            xlabel={time $t$ in \si{\day}},
            width=0.5\linewidth,
            height=0.4\linewidth,
            legend columns = 4,
            legend style = {anchor=north, style={column sep=3pt}},
            grid=both,
            grid style={line width=.1pt, draw=gray!10},
            y tick label style={
                /pgf/number format/fixed zerofill,
                /pgf/number format/precision=1},
            scaled ticks=false, tick label style={/pgf/number format/fixed}
        ]
        \nextgroupplot[
            xmin=0, xmax=29,
            ymin=0, ymax=1.250,
            ytick={0,.250,...,1.250},
            restrict x to domain=0:42,
            xtick={1,9,21,28},
            axis y line*=left,
            y tick label style={
                /pgf/number format/fixed zerofill,
                /pgf/number format/precision=2},
        ]
            \addplot+ [draw=none, no markers, name path=A,forget plot] table [skip first n=3, x index=1, y index=5, col sep=comma, filter discard warning=false] {figures/hydrate-rim/data_t_suess/CSH-length_C3S-C2S.csv};
            \addplot+ [draw=none, no markers, name path=B,forget plot] table [skip first n=3, x index=1, y index=7, col sep=comma, filter discard warning=false] {figures/hydrate-rim/data_t_suess/CSH-length_C3S-C2S.csv};
            \addplot [fill=red, fill opacity=0.2,forget plot] fill between [of=A and B];

            \addplot [color=red, mark=x, dotted] table [skip first n=3, x index=1, y index=6, col sep=comma] {figures/hydrate-rim/data_t_suess/CSH-length_C3S-C2S.csv};


            \addplot+ [draw=none, no markers, name path=A,forget plot] table [skip first n=3, x index=1, y index=14, col sep=comma, filter discard warning=false] {figures/hydrate-rim/data_t_suess/CSH-length_C3S-C2S.csv};
            \addplot+ [draw=none, no markers, name path=B,forget plot] table [skip first n=3, x index=1, y index=16, col sep=comma, filter discard warning=false] {figures/hydrate-rim/data_t_suess/CSH-length_C3S-C2S.csv};
            \addplot [fill=green, fill opacity=0.2,forget plot] fill between [of=A and B];

            \addplot [color=green, mark=x, dotted] table [skip first n=3, x index=1, y index=15, col sep=comma, filter discard warning=false] {figures/hydrate-rim/data_t_suess/CSH-length_C3S-C2S.csv};


            \addplot+ [draw=none, no markers, name path=A,forget plot] table [skip first n=3, x index=0, y index=1, col sep=comma, filter discard warning=false] {figures/hydrate-rim/data_t_suess/CSH-length_C2S_moeser.csv};
            \addplot+ [draw=none, no markers, name path=B,forget plot] table [skip first n=3, x index=0, y index=3, col sep=comma, filter discard warning=false] {figures/hydrate-rim/data_t_suess/CSH-length_C2S_moeser.csv};
            \addplot [fill=cyan, fill opacity=0.2,forget plot] fill between [of=A and B];

            \addplot [color=cyan, mark=x, dotted] table [skip first n=3, x index=0, y index=2, col sep=comma, filter discard warning=false] {figures/hydrate-rim/data_t_suess/CSH-length_C2S_moeser.csv};

        \nextgroupplot[
            xmin=-65, xmax=1000,
            restrict x to domain=28:1000,
            xtick={28,150, 360, 600, 1000},
            ytick={.5,1,1.5,...,2.5},
            axis y line*=right,
            axis x discontinuity=parallel,
            legend to name={CommonLegend2}
        ]
            \pgfplotsset{cycle list shift=+1}
            \addplot+ [draw=none, no markers, name path=A,forget plot] table [skip first n=3, x index=1, y index=5, col sep=comma, filter discard warning=false] {figures/hydrate-rim/data_t_suess/CSH-length_C3S-C2S.csv};
            \addplot+ [draw=none, no markers, name path=B,forget plot] table [skip first n=3, x index=1, y index=7, col sep=comma, filter discard warning=false] {figures/hydrate-rim/data_t_suess/CSH-length_C3S-C2S.csv};
            \addplot [fill=red, fill opacity=0.2,forget plot] fill between [of=A and B];

            \addplot [color=red, mark=x, dotted] table [skip first n=3, x index=1, y index=6, col sep=comma, filter discard warning=false] {figures/hydrate-rim/data_t_suess/CSH-length_C3S-C2S.csv};
            \addlegendentry[]{\ac{C3S} \cite{voland.2000,suss.2001,stark.2001,moser.2002}}


            \addplot+ [draw=none, no markers, name path=A,forget plot] table [skip first n=3, x index=1, y index=14, col sep=comma, filter discard warning=false] {figures/hydrate-rim/data_t_suess/CSH-length_C3S-C2S.csv};
            \addplot+ [draw=none, no markers, name path=B,forget plot] table [skip first n=3, x index=1, y index=16, col sep=comma, filter discard warning=false] {figures/hydrate-rim/data_t_suess/CSH-length_C3S-C2S.csv};
            \addplot [fill=green, fill opacity=0.2,forget plot] fill between [of=A and B];

            \addplot [color=green, mark=x, dotted] table [skip first n=3, x index=1, y index=15, col sep=comma, filter discard warning=false] {figures/hydrate-rim/data_t_suess/CSH-length_C3S-C2S.csv};
            \addlegendentry[]{$\beta$-\ac{C2S} \cite{voland.2000,suss.2001}}


            \addplot+ [draw=none, no markers, name path=A,forget plot] table [skip first n=3, x index=0, y index=1, col sep=comma, filter discard warning=false] {figures/hydrate-rim/data_t_suess/CSH-length_C2S_moeser.csv};
            \addplot+ [draw=none, no markers, name path=B,forget plot] table [skip first n=3, x index=0, y index=3, col sep=comma, filter discard warning=false] {figures/hydrate-rim/data_t_suess/CSH-length_C2S_moeser.csv};
            \addplot [fill=cyan, fill opacity=0.2,forget plot] fill between [of=A and B];

            \addplot [color=cyan, mark=x, dotted] table [skip first n=3, x index=0, y index=2, col sep=comma, filter discard warning=false] {figures/hydrate-rim/data_t_suess/CSH-length_C2S_moeser.csv};
            \addlegendentry[]{$\beta$-\ac{C2S} \cite{moser.2002}}

        \end{groupplot}
        \node [above] (L) at (rel axis cs:1.0,1.14) {\ref{CommonLegend2}};

    \end{tikzpicture}
    \caption{Growth of \ac{CSHo} on $\beta$-\ac{C2S} compared to \ac{C3S} at \SI{20}{\celsius}. Based on data from \textcite{voland.2000}, \textcite{suss.2001}, \textsc{Stark} and \textsc{Möser} et al. \cite{stark.2001,moser.2002}.}
    \label{fig:suess_CSH_C2S}
\end{figure}

In his presentations, \textsc{Möser} summarized the results as follows:
\begin{iquote}
When growing unhindered, the needle-like \ac{CSH} phases can reach a length of approximately \SI{1.4}{\micro\meter}. This applies only for the hydration of alite. In pure belite, a \ac{CSH}-length of up to \SI{2.5}{\micro\meter} can be measured. \cite{moser.2001}
\end{iquote}
%Zitat Möser
%It can be assumed that at the beginning of the acceleration period all C-S-H phases grow in form of fibres. These fibres grow in length during the course of hydration with hardly any changes in thickness. When growing unhindered, the needle-like \ac{CSH} phases can reach a length of approx. 1.4 µm. These applies only for hydration of alite (in pure belite: \ac{CSH}- length up to 2.5 µm). In compact concrete it is to be expected that there is intergrowth between the fibers as well as and interlocking with one another and thus gel pores are formed which have often been described.


\subsection{Hadley grains}
\label{subsec:hadley}

Sometimes, in polished sections of hydrated cementitious materials, hollow hydrate rims, sometimes filled with an unhydrated clinker grain can be observed (see Figure~\ref{fig:hadley}).
These objects are called \textsc{Hadley grains} after \textcite{hadley.1972} who first described this phenomenon \cite{hadley.2000} in detail.
\textcite{kjellsen.1997} described the formation of hollow shells depending on the \ac{wb}-ratio and silica fume content.
They saw a increasing porisity attributed to \textsc{Hadley grains} with an increasing \ac{wb}-ratio especially for young specimens.
\ac{PXCT} enabled \textcite{shirani.2023} to analyze the early development of hadley grains around \ac{C3S} grains.
They identified a gap of \SI{450}{\nano\meter} between \ac{CSH} and the unhdydrated \ac{C3S}.
This \ac{CSH} had a very low density compared to more matured material.

\textcite{yio.2015a} observed the development of such a grain using confocal microscopy.
After \SI{28}{\day} they observed `\textit{needle-like hydration products that connect the outer products and the anhydrous core}'\cite{yio.2015a}.
In the later stages, the grain can completely dissolve and empty shells of \ac{CSH} remain.
This can be seen in Figure~\ref{fig:hadley}, where both filled and hollow Hadley grains are visible, especially around belite.

\begin{figure}[htb]
	\centerline{\includegraphics[width=\textwidth]{hydrate-rim/hadley_grains.pdf}}
	\caption{BSE image of a polished section of a \SI{28}{\day} hydrated alite and belite mixture (3:1). Around many belite grains, Hadley grains were formed. Furthermore, empty \ac{HLT} shells can be found. \label{fig:hadley}}
\end{figure}